 \documentclass{article}
\usepackage{tikz}
\usetikzlibrary{angles, quotes}
\usepackage{amssymb}
\usepackage{pgfplots}
\pgfplotsset{width=10cm,compat=1.9}
\usepgfplotslibrary{statistics}
\usepackage[left=2.54cm,right=2.54cm,top=2.54cm,bottom=2.54cm]{geometry}
\usepackage{tasks} %Paquete necesario para la producción de items horizontales para algunos ejercicios de tarea empleando el comando /begin{multicols}}{#}.../end{multicols} para ayudarnos a reducir las páginas
\usepackage{pdflscape} %comando que permite cambiar a orientación horizontal las hojas%
\usepackage{soul} %Paquete para permitir la compilación de acentos usando el comando {\hl{}}}\\
\sethlcolor{yellow(munsell)} %comando para definir el color para el subrayado usando el comando \hl
\usepackage[utf8]{inputenc}
\usepackage[spanish]{babel}
\usepackage{icomma}
\usepackage{ragged2e}
\usepackage{siunitx}
\usepackage{url}
\usepackage[colorlinks=true, urlcolor=blue,  linkcolor=black, citecolor=green]{hyperref}
\usepackage{pdfpages}
\usepackage{blindtext} %Paquete que genera texto ficticio (dummy text) con el comando: \blindtext
\usepackage{cancel} %in the preamble gives you four different modes of striking through: \cancel{text to cancel} draws a diagonal line (slash) through its argument, \bcancel{text to cancel} uses the negative slope (a backslash), \xcancel{text to cancel} draws an X (actually \cancel plus \bcancel), \cancelto{〈value〉}{〈expression〉} draws a diagonal arrow through the 〈expression〉pointing to the 〈value〉 (math-mode only)
\usepackage{csquotes}
\usepackage{afterpage}
\usepackage{parskip} 
\usepackage{float}
\usepackage{enumitem}
\usepackage{multicol}%Paquete que permite la creación aislada de columnas de texto. De esta forma se reduce la cantidad de páginas en nuestro documento
\usepackage{lipsum} %Paquete que genera texto de relleno al igual que \usepackage{blindtext}, se usa con el comando \lipsum[#-#] los #(hashtags) delimitan el número de páginas que deseamos ocupar del paquete lipsum para usarlos como relleno. 

% Definir el comando personalizado para la nota
\newcommand{\mynote}[1]{%
\begin{tikzpicture}[baseline=-0.75ex]
\node[draw, circle, fill=Apple Green!60, inner sep=2pt] (note) {\textbf{Nota:}};
\end{tikzpicture}%
\ \textit{#1}%
}

\newenvironment{Figura}
{\par\medskip\noindent\minipage{\linewidth}}
{\endminipage\par\medskip}
\usepackage{caption}
\usepackage[
backend=biber,
sorting=none,
url=true
]{biblatex} %bibliografía
\addbibresource{biblio.bib}
% Establecer el espacio entre las entradas de la bibliografía
\setlength{\bibitemsep}{1\baselineskip} % Puedes ajustar el valor según tus preferencias
\usepackage{amsmath,amsthm,amssymb,amsfonts}
\usepackage{pifont} %Permite la compilación del comando \xmark: es el símbolo de la tachita
\usepackage{mathtools} %permite la compilación de símbolos para matrices
\usepackage{empheq} %Paquete que se relaciona con \usepackage[most]{tcolorbox} para la creación de cuadros/cajas de colores para delimitar los resultados matemáticos o líneas de texto usando la línea de comando: \begin{empheq}[box={\mymath[colback=orange(webcolor)!70,drop lifted shadow]}]{equation*} \end{empheq}
\usepackage[most]{tcolorbox} %Paquete que se relaciona con \usepackage{empheq} para la creación de cuadros/cajas de colores para delimitar los resultados matemáticos o líneas de texto
\renewcommand{\qedsymbol}{$\blacksquare$}
\usetikzlibrary{positioning,decorations.pathreplacing} %Paquete que permite la compilación de llaves para cuadros sinópticos (brace diagram) 
\usepackage{schemata} %The schemata package is designed just for these brace diagrams
\usepackage{booktabs}

\newcommand\AB[2]{\schema{\schemabox{#1}}{\schemabox{#2}}} %comando o paquete necesario para crear cuadros sinópticos

\newtcbox{\mymath}[1][]{%
nobeforeafter, math upper, tcbox raise base,
enhanced, colframe=blue!30!black,
colback=blue!30, boxrule=1pt,
#1} %Comando importante para encasillar los resultados matemáticos en cajas de diferentes colores, se relaciona con los paquetes \usepackage{empheq} y \usepackage[most]{tcolorbox} 

\newenvironment{sysmatrix}[1]
{\left(\begin{array}{@{}#1@{}}}
{\end{array}\right)}
\newcommand{\ro}[1]{%
\xrightarrow{\mathmakebox[\rowidth]{#1}}%
}
\newlength{\rowidth}% row operation width
\AtBeginDocument{\setlength{\rowidth}{3em}}  %Comando importante para laproducción de lineas de operaciones entre matrices, método de Gauss o Gauss Jordan se relaciona con \newenvironment{sysmatrix}

%formato para cambiar el horario a español
\usepackage[spanish]{datetime2}
\DTMsetdatestyle{spanish}

\renewcommand{\today}{\DTMdisplaydate{\the\year}{\the\month}{\the\day }{-1}}
%%%%%%%%%%%%%%%%%%%%%%%%%%%%%%%%%%%%%%%%%

\usepackage{datetime}
\newcommand{\mycurrenttime}{\xxivtime}

%%%%%%%%%%%%%%%%%%%%%%%%%%%%%%%%%%%%%%%%%%%%%%%%%%%%%%%%%%%%%%%%%%%%%%%%%%%%%%%%%%%%%%%%%%%%%%%%%%%%%%%%%%%%%%%%%%%%%%%%%%%%%%%%%%%%%%%%%%%%%%%%%


%%%%%%%%%%%%%Caja de color para el título%%%%%%%%%%%%%%%%%%%%%%%%%%%%%%%%%%%%%%%%%%%%%%%%%%%%%%

\definecolor{myframecolor}{RGB}{1, 45, 75} % Prussian Blue
\definecolor{myboxcolor}{RGB}{0, 107, 92} % Apple Green

\newtcolorbox{mybox}{
enhanced,
colback=myboxcolor!7, % Color de fondo del cuadro
colframe=myframecolor, % Color del marco del cuadro
arc=0pt,
boxrule=1pt,
borderline west={2mm}{-10mm}{myframecolor}, % Borde en el lado izquierdo
sharp corners=southwest,
width=\linewidth,
before=\par\vspace{\bigskipamount}, % Espacio antes del cuadro
after=\par\vspace{\bigskipamount} % Espacio después del cuadro
}


\newcommand{\euler}{e} %Comando para producir letra e de euler
\newenvironment{remark}{\par\vfill\footnotesize % Vertical white space above the remark and smaller font size
\begin{list}{}{
	\leftmargin=80pt % Indentation on the left
	\rightmargin=60pt}\item\ignorespaces % Indentation on the right
\makebox[-2.5pt]{\begin{tikzpicture}[overlay]
		\node[draw=Horizon!60,line width=2.5pt,circle,fill=Horizon!25,font=\sffamily\bfseries,inner sep=4pt,outer sep=0pt] at (-19pt,5pt){\textcolor{Horizon}{Nota}};\end{tikzpicture}} % Blue Nota in a circle
\advance\baselineskip -1pt}{\end{list}\vskip5pt} % Tighter line spacing and white space after remark
\usepackage{graphicx}

\usepackage{titling}

\input{colores.tex}

%%Producir signos de cita textual``Asignatura''

\newcommand{\xmark}{\ding{55}} %%%comando que genera la tachita

\newenvironment{MyColorPar}[1]{%
\leavevmode\color{#1}\ignorespaces%
}{%
}%

%%%%%%%%%%%%%%%%%%% Título de la tarea, Nombre de alumno

\title{Clase \#31 - Electrodinámica}
\author{\emph{JA} $\boldsymbol{\mid}$}
\date{\today}

%%%%%%%%%%%%%%%%%%% Datos de la Materia

\usepackage{fancyhdr}
\fancypagestyle{plain}{%  the preset of fancyhdr 
\fancyhf{} % clear all header and footer fields
\fancyfoot[R]{
\includegraphics[width=2cm]{LogoFCFMBUAP (1).png}}
\fancyfoot[L]{{\bfseries{\thedate{}}} a las {\bfseries{\mycurrenttime{}}} horas{} (GMT-6, H. Puebla de Zaragoza, Pue)}
\fancyhead[L]{Electrodinámica}
\fancyhead[R]{\theauthor}
}
\makeatletter
\def\@maketitle{%
\newpage
\null
\vskip 1em%
\begin{center}%
\let \footnote \thanks
{\LARGE \@title \par}%
\vskip 1em%
%{\large \@date}%
\end{center}%
\par
\vskip 1em}
\makeatother

\usepackage{lipsum}  
\usepackage{cmbright}


%%%%%%%%%%%%%%%%%%%%%%%%%%%%%%%%%%%%%%%%%%%%%%%%%%%%%%%%%%%%%%%%%%%%%%%%%%%%%%%%%%%%%%%%%%%%%%%%%%%%%%%%%%%%%%%%%%%%%%%%%%%%%%%%%%%%%%%%%%%%%%%%%%%%%%%%%%%%%%%%%%%%%%%%%%%%%%%%%%%%%%%%%%%%%%%%%%%%%%

\begin{document}



\maketitle


%%%%%%%%%%%%%%%%%%% Datos del profesor, campus y aula

\begin{mybox}
\noindent\begin{tabular}{@{}ll}
	{\bfseries{Profesor}} & Dr. Iván Fuentecilla Cárcamo\\
	\textcolor{prussianblue}{\bfseries{Campus}} / \textcolor{trueblue}{\bfseries{Facultad}}  & \textcolor{prussianblue}{\bfseries{C.U. BUAP}} /  \textcolor{trueblue}{\bfseries{FCFM}} \\
	\textcolor{prussianblue}{\bfseries{Edificio}} /  \textcolor{trueblue}{\bfseries{Salón}}    & 	\textcolor{prussianblue}{\bfseries{EMA4}} /  \textcolor{trueblue}{\bfseries{409}}   
\end{tabular} 
\end{mybox} \vspace{1cm}




	





\begin{MyColorPar}{orange(colorwheel)}

\end{MyColorPar}


\tableofcontents \newpage

\begin{center}
\begin{tikzpicture}
	% Dibuja círculos rellenos en diferentes colores
	\fill[bittersweet] (0,0) circle (0.17cm);
	\fill[Apple Green] (0.6,0) circle (0.17cm);
	\fill[trueblue] (1.2,0) circle (0.17cm);
	\fill[Aguamarina] (1.8,0) circle (0.17cm);
	\fill[Sun] (2.4,0) circle (0.17cm);
	\fill[orange(colorwheel)] (3.0,0) circle (0.17cm);
	\fill[Pesto] (3.6,0) circle (0.17cm);
	\fill[blue-violet] (4.2,0) circle (0.17cm);
	\fill[ticklemepink] (4.8,0) circle (0.17cm);
	\fill[Gallery] (5.4,0) circle (0.17cm);
	\fill[Mantle] (6.0,0) circle (0.17cm);
	\fill[Nevada] (6.6,0) circle (0.17cm);
	\fill[onyx] (7.2,0) circle (0.17cm);
	% Agrega más círculos o modifica los colores según sea necesario
\end{tikzpicture}
\end{center}  \vspace{2cm}




\section{Unit 2: Electric fields in dielectric media}

Tenemos:


\[ \epsilon, \, k,\, \mathcal{x} \]


 \begin{figure}[H]
    \centering
	    \includegraphics[width=0.7\textwidth]{1.png} % ajusta el ancho según convenga
    \caption{Descripción clara y concisa de la figura.}
    \label{fig:mi_figura} % opcional: para usar con \ref{fig:mi_figura}
\end{figure}


\[ \vec{P} = \lim\limits_{\Delta{V} \to 0} \frac{ \Delta{\vec{P}} }{\Delta{V}} \]


 \begin{figure}[H]
    \centering
	    \includegraphics[width=0.7\textwidth]{2.png} % ajusta el ancho según convenga
    \caption{Descripción clara y concisa de la figura.}
    \label{fig:mi_figura} % opcional: para usar con \ref{fig:mi_figura}
\end{figure}


Tenemos:


\begin{equation}
\begin{split}
 \vec{E}_{m} &= \vec{E}_{x} + \vec{E}_{P} + \vec{E}_{S} + \vec{E}^{\prime}\\
 \sigma_{P} &= \vec{P} \cdot \hat{n}\\
 \rho_{P}     &= -\nabla \cdot \vec{P}
\end{split}
\end{equation}

 \begin{figure}[H]
    \centering
	    \includegraphics[width=0.7\textwidth]{3.png} % ajusta el ancho según convenga
    \caption{Descripción clara y concisa de la figura.}
    \label{fig:mi_figura} % opcional: para usar con \ref{fig:mi_figura}
\end{figure}

\[ \sigma_{p} = \vec{P} \cdot \hat{r} = P\cos \theta  \]


\[ \boxed{  \vec{E} = \frac{ 1}{  4\pi \epsilon_{0}  } \frac{ q }{ \left| \vec{r} - \vec{r}^{\prime} \right|^{3} } ( \vec{r} - \vec{r}^{\prime}  )    } \]

$P\cos\theta$ es el elemento de carga, la $\vec{r}$ es el punto de observación.

Tenemos $\ldots$

\[ \vec{E}_{m} = \vec{E} + \vec{E}_{s} + \vec{E}^{\prime}     \]

\[ d\vec{E}_{s} = \frac{ -P\cos \theta \vec{r} da    }{4\pi \epsilon_{0} r^{3} }  \]

\textbf{Vector unitario en coordenadas esféricas}


\[ \hat{r} = \cos \phi \sen \theta \hat{i} + \sen \phi \sen \theta \hat{j} + \cos \theta \hat{k} \]


 \begin{figure}[H]
    \centering
	    \includegraphics[width=0.7\textwidth]{4.png} % ajusta el ancho según convenga
    \caption{Descripción clara y concisa de la figura.}
    \label{fig:mi_figura} % opcional: para usar con \ref{fig:mi_figura}
\end{figure}

La polarización apunta en la dirección de $z$. Entonces podremos cancelar términos.

De estas componente que son 3, solo sobrevive una. Es justamente $dE_{zs}$
Escribir lo que va a resultar después de hacer esa consideración, con $da = r^{2}\sen\theta d\theta d\phi$


Tenemos:



\begin{equation}
\begin{split}
\int\limits E_{sz} &= \frac{-P \cos \theta \vec{r}da}{4\pi \epsilon_{0} r^{3}} \quad \text{tomando} \boxed{da= r^{2}\sen\theta d\theta d\phi}\\
 \int\limits dE_{sZ}  &=\int\limits \frac{ -P \cos \theta \hat{k} r^{2}\sen\theta d\theta d\phi   }{ 4\pi \epsilon_{0}r^{3}}\\
 E_{sz}  &=  \frac{-P}{4\pi \epsilon_{0}r} \hat{k}  \int\limits \cos\theta \sen\theta d\theta d\phi 
\end{split}
\end{equation}

Tenemos bien realizado:

Considerando que $\hat{i}$ y $\hat{j}$ se van a cero.
 
 Solo queda la componente en $z$ del vector unitario $r$ en coordenadas esféricas, tenemos así:
 
 
 \begin{equation}
 \begin{split}
  dE_{z}  &=  \frac{   -P \cos \theta r \cos \theta r^{2}\sen \theta d\theta d\phi    }{4\pi \epsilon_{0}r^{3}  }\\
  dE_{z} &=  \frac{ P \cos^{2}\theta \sen \theta d\theta d\phi }{4\pi \epsilon_{0}}\\
   d\vec{E}  &= \vec{P} \int\limits_{0}^{2\pi} \int\limits_{0}^{\pi} \frac{1 }{4\pi \epsilon_{0}} \cos^{2}\theta \sen \theta d\theta d\phi,\quad  \text{tenemos},\quad \int\limits_{0}^{\pi} \cos^{2} \theta \sen\theta d\theta = \left.- \frac{\cos^{3} \theta}{3}\right|_{0}^{\pi}  = \boxed{ \frac{2}{3}}\\
  \vec{E} &= \frac{ \vec{P} }{3\epsilon_{0}} 
 \end{split}
 \end{equation}



Tenemos: 


\[ \boxed{ \vec{E}_{s} = \frac{ \vec{P}}{3\epsilon_{0}}} \]

\begin{equation}
\begin{split}
\vec{P}_{m}   &\underbrace{\propto}_{Polarizability} \vec{E}_{m} \\
\vec{P} &= \mathcal{N} \propto \vec{E}_{m}\\
\vec{P} &= \mathcal{N} \propto \left(  \vec{E} + \frac{ \vec{P}  }{3\epsilon_{0}}    \right)
\end{split}
\end{equation}

El campo molecular está generando un momento dipolar en la molécula.

Tenemos $\ldots$


\begin{equation}
\begin{split}  
 \vec{P}    &= \mathcal{N} \propto \left(  \vec{E} + \frac{\vec{P}  }{ 3\epsilon_{0} }  \right)\\
 \vec{P} &= \mathcal{X} \vec{E}\\
 D &= \epsilon \vec{E} = \epsilon_{0}\vec{E}+ \vec{P} = \epsilon_{0}\vec{E}+\mathcal{X}\vec{E}\\
  \mathcal{K} &= \frac{\epsilon}{\epsilon_{0}}
\end{split}
\end{equation}

$\vec{P}$ es la polarización macroscópica $P$.


Tenemos que hacer de actividad; \textbf{vincular lo macroscopico con lo microscopico}.  A través d ela medición de la constante dieléctrica con lo molecular, la polarizabilidad molecular. La $\alpha$ nos da los efectos de la polarizabilidad.


Calcular:

\[ \boxed{ \alpha( \mathcal{K} )   } =  \]

Tenemos: 

\begin{equation}
\begin{split}
 \vec{P} &= \mathcal{N} \alpha \left(   \vec{E} + \frac{ \vec{P} }{3 \epsilon_{0}} \right)\\
\mathcal{X}\vec{E}   &= \mathcal{N} \alpha \left(  \vec{E} + \frac{ \mathcal{X}\vec{E} }{3\epsilon_{0}} \right)\\
\mathcal{X}\vec{E} &	= \mathcal{N} \alpha \left(    \frac{3\epsilon_{0}\vec{E} + \mathcal{X}\vec{E}}{ 3\epsilon_{0}  } \right) \\
\mathcal{X}\vec{E} &= \frac{1}{3\epsilon_{0}}\mathcal{N}\alpha \left( 3\epsilon_{0}\vec{E} + \mathcal{X}\vec{E}\right)\\
3\epsilon_{0}\mathcal{X}\vec{E} &= \mathcal{N}\alpha \left( 3\epsilon_{0}\vec{E} + \mathcal{X}\vec{E}\right)\\
3\epsilon_{0} \mathcal{X} &= \mathcal{N} \alpha \left( 3\epsilon_{0} + \mathcal{X} \right)\\
3\epsilon_{0} (k-1)\epsilon_{0} &= \mathcal{N}\alpha (3\epsilon_{0} + (k-1)\epsilon_{0})\\
3\epsilon_{0} (k-1) &= \mathcal{N}\alpha (3 + (k-1))\\
3\epsilon_{0} (k-1) &= \mathcal{N}\alpha (2 + k)\\
\frac{3\epsilon_{0} (k-1)}{\mathcal{N} (k+2)}  &= \alpha 
\end{split}
\end{equation}

Al resultado anterior se le conoce como:

$\boxed{ \text{Clausius - Mossoti equation } }$

 \begin{figure}[H]
    \centering
	    \includegraphics[width=0.7\textwidth]{5.png} % ajusta el ancho según convenga
    \caption{Descripción clara y concisa de la figura.}
    \label{fig:mi_figura} % opcional: para usar con \ref{fig:mi_figura}
\end{figure}

 \begin{figure}[H]
    \centering
	    \includegraphics[width=0.7\textwidth]{5.png} % ajusta el ancho según convenga
    \caption{Descripción clara y concisa de la figura.}
    \label{fig:mi_figura} % opcional: para usar con \ref{fig:mi_figura}
\end{figure}

Tenemos que irnos a la sección de Reit Mildford 5.2 dipolos inducidos, un modelo simple.

deducir paso a paso ecuación 5-13

\end{document}